\chapter*{前言}
\addcontentsline{toc}{chapter}{前言}
\markboth{前言}{前言}

从公元前206年西汉王朝建立，到公元581年隋文帝统一中国，历时近八百年，史称两汉魏晋南北朝时期，历史上也有人称这一时期为汉魏六朝。这一时期的文学在继承先秦文学发展的基础上，开创了文学发展的新局面，出现了许多新文体、新思路、新风貌，特别是诗文，获得了更加长足的发展。应当承认这一时期的文学是空前繁荣的，其成就也是辉煌的。也正因为如此，汉唐文化才能成为历史上中华民族先进文化的代表。文学的其他方面的发展与成就我们姑且不论，现在仅就汉魏六朝文的发展与成就略抒己见。

汉代的文引人注意的有四个板块，一是西汉初以贾谊、晁错为代表的政论文，本书选注了贾谊的《过秦论》和晁错的《论贵粟疏》作为代表。这一类文章的出现自有它的时代背景的：刘邦建立了西汉王朝之后，为了求得王朝的长治久安，开始注意总结历史上王朝兴败盛衰的原因，特别是秦朝速亡的原因，更是他的前车之鉴，《过秦论》、《论治安策》等，便是在这一背景下产生的。同时，汉初的经济在经过了秦末的社会大动乱之后，本来就不太发达的生产力，又遭到了严重的破坏，当时的社会还潜伏着许多的矛盾和危机，这种状况直到汉文帝时代，还不见太大的转机，所以贾谊在《陈政事疏》说：

\begin{quote}
臣窃惟事势，可为痛哭者一，可为流涕者二，可为长太息者六，若其它背理而伤道者，难遍以疏举。进言者皆曰天下已安已治矣，臣独以为未也。曰安且治者，非愚则谀，皆非事实知治乱之体者也。夫抱火厝之积薪之下而寝其上，火未及燃，因谓之安，方今之势，何以异此！
\end{quote}

贾谊是没有胆量在皇帝面前危言耸听的，他不过是敢于正视现实又敢于讲真话罢了。晁错也与贾谊有点相似，他也是个敢讲真话的人，为了加强中央集权，他极力主张削弱诸侯王的势力，为此他付出了血的代价。他的《论贵粟疏》也能洞察当务之急，所言都颇中肯綮。鲁迅先生在《汉文学史纲要》中说贾谊、晁错的文章“疏直激切”，其评价是中肯的。

汉文的第二个板块是汉赋，它是整个汉代最主要的文体，是纯文学作品，也是汉代发展演变最为迅速的文体。刘勰的《文心雕龙·诠赋》篇说赋是“受命于诗人，拓宇于楚辞”。说它“受命与诗人”，其着眼点是《诗大序》有“风、赋、比、兴、雅、颂”的“六义”之说，赋的得名即从“六义”之中位居其二的“赋”而来。“六义”中的赋，不过是一种表现手法，并不是一种文体。从汉赋开始，变成了一种文体，而这种文体的特点便是“铺采摛文，体物写志”（《文心雕龙·诠赋》）。所谓“拓宇于楚辞”，就是说汉赋是从《楚辞》的领域开拓发展而来。从赋的发展史而言，楚辞是骚赋的时代，骚赋的写作一直延续到汉初，王逸的《楚辞章句》就收有汉人的作品。西汉人的骚赋在体式上与《楚辞》是十分相似的，如我们所选的贾谊的《鵩鸟赋》与《吊屈原赋》，在句式上还保留着《楚辞》的“兮”字，把它们视为骚赋也并没有什么不妥。可是枚乘的《七发》就与贾谊的赋有所不同了，不仅没有了“兮”字，在描写上铺张扬厉之风也开始形成，赋中也有了主客的对话，他虽然没有以赋名篇，实际上在体式上已接近大赋，可视为骚赋向大赋（辞赋）的过渡。萧统的《文选》没有把《七发》放在赋类里，而是单设了一类“七”体，我想萧统所以这样做，可能有两个考虑。一是《七发》的体式确实与赋有所不同，《七发》虽有吴客和楚太子的对话，但并没有主客之间或客与客之间的互相辩难。在《七发》中楚太子没有说几句话，吴客却连用七事来循循善诱地启发太子，可以说《七发》就是启发，就是用七件事来启发太子，最后以太子据几而起曰：“听圣人辩士之言，涊然汗出，霍然病已”作结，这与汉代大赋以“卒章显其志”、“曲终奏雅”的方式最后来一点讽谏又有所不同。从严格的意义上说，《七发》还不能算赋，此其一。其二是，枚乘《七发》问世之后，仿效者颇多，逐渐形成了一种体式，这就是《七》体。曹植在其所写的《七启·序》中说：

\begin{quote}
昔枚乘作《七发》，傅毅作《七激》，张衡作《七辩》，崔骃作《七依》，辞各美丽，余有慕之焉。遂作《七启》。并命王粲作焉。
\end{quote}

这说明在曹植的时代，《七》体已初具规模。《文选》把它列为一类，也自有它的道理。另外《七发》还不能说是“拓宇于楚辞”的，它的体式与风格与楚辞没有多少相似之处，倒是与战国时代的游说之士的言谈风格有些相似，所以楚太子把吴客之所言，视为“圣人辩士之言”。但《七发》在赋的发展史上也应占有一席地位，它是向赋过渡的桥梁。

汉代的大赋的模式是司马相如奠定的，我们所选注的司马相如的《子虚》、《上林》，实则是《天子游猎赋》一篇大赋的两个组成部分。此赋在艺术描写上极尽铺张扬厉之能事，赋中又虚构了三个人物：子虚、乌有、亡是公，有主有客，有互相之间的对话，也有互相之间的辩难。楚国使者子虚先生在看到齐王举行的盛大的游猎活动之后，在齐国的乌有先生与天子使者亡是公面前大夸楚国的园囿之大、游猎之盛、物产之富，以期压倒齐国。乌有先生加以辩难。最后天子的使者亡是公出来，盛夸天子的上林苑之大、车骑游猎之壮观，皆非齐楚可比。又通过亡是公之口，说出天子已觉察到这样做未免太奢侈了，“于是乎乃解酒罢猎，而命有司，曰：‘地（指上林苑）可垦辟，悉为农郊，以赡萌隶。隤墙填堑，使山泽之人得至焉。……发仓廪以救贫穷，补不足，恤鳏寡，存孤独。出德号，省刑罚，改制度，易服色，革正朔，与天下为更始。’”这就是“其卒章归之于节俭。因以风谏”（《史记·司马相如列传》）。大赋的“体物写志”也于此可见。

自司马相如奠定了汉赋的模式之后，两汉的大赋模式基本上没有太大的变化，但汉代赋的发展并没有凝固化，而是处在迅速的变化之中，这不仅表现在题材的开拓上，也表现在体式的发展变化上。

《文选》六十卷，从卷一到卷十九都是选的赋，几乎占了三分之一的篇幅。它把赋又分为十五类，这是从题材内容上分类的。有些分类在今天看来也并不完全科学，但也初步体现出汉赋内容的丰富性，特别是描写京都与畋猎的大赋，闳丽巨衍，雍容典雅，它是汉代国力富强的表现，是汉代文人润色鸿业的力作，也是当时的时代产物。要正确而全面的评价汉赋，不能不看到这些方面。

对于汉赋如何评价，在汉代就有所争论，司马迁对司马相如的赋已略露微词，他指出相如曾向汉武帝献《大人赋》，本意是讽刺汉武帝好神仙之术的，汉武帝看了之后反而飘飘然有凌云之意，可见赋的讽刺效果是很不理想的。司马迁又说：“相如虽多虚辞滥说，然其要归引之节俭，此与《诗》之风谏何异。”（《史记·司马相如列传》）既批评了司马相如赋的“虚辞滥说”又肯定了它的讽谏作用，其评价还算公允。而扬雄对汉赋的批评，比司马迁更为激烈。他在《法言·吾子》篇中，把赋看作是“童子雕虫篆刻”的玩意，是“壮夫不为”的。他本身就是赋家，却对赋如此的轻视，令人难以理解。当谈到赋的讽刺作用时，他认为赋的讽刺作用是劝百讽一的：“讽则已，不已，吾恐不免于劝也。”汉代的文论家比较重视文学的教化作用，对文学的审美作用则不太重视，所以批评汉赋的虚辞滥说、雕虫篆刻还是可以理解的，有它合理的因素。司马相如奠定的汉赋模式，自有它自己的局限，所以在东汉后期，出现了张衡的《归田赋》、赵壹的《刺世疾邪赋》、王粲的《登楼赋》等抒情小赋，它们汰除了铺张扬厉之风，以及堆垛寡变、板滞凝重之弊，用清新俊永、短小精悍的抒情小赋，取代了闳丽巨衍的大赋，也增强了赋的抒情性，是赋史上的一次变迁。

两汉文的第三个板块是以《史记》、《汉书》为代表的历史纪传文，因限于篇幅，本书没有选录。鲁迅先生称《史记》为“史家之绝唱，无韵之《离骚》”（《汉文学史纲要》）。其思想与艺术成就均在班固的《汉书》之上，在汉文中占有重要的地位。

两汉文的第四个板块是应用文，应用文的范围很广，章、表、书、记、碑、诔等都包括在内。汉代的应用文起点很高，它已发展到相当高的水平，像司马迁的《报任少卿书》，可以说是用血泪写成的，今天读来，仍然令人回肠荡气，叹为观止。

汉代写碑文的以蔡邕最为有名，刘勰《文心雕龙·诔碑》篇云：“自后汉以来，碑碣云起，才锋所断，莫高蔡邕。”我们只选了蔡邕的《郭有道碑文》一篇。东汉的“谀墓”之风颇为盛行，有钱的人家只要肯出高价，便可以对碑主大加美化，以至出现“生为盗跖，死为夷齐”的现象。蔡邕曾对人说：“吾为碑铭多矣，皆有惭德，唯《郭有道》无愧色耳。”（《后汉书·郭太传》）可见这篇碑文是比较真实的。

建安时代孔融和陈琳的应用文成就也是比较高的。孔融有两篇举荐人物的文章，一是《论盛孝章书》，一是《荐祢衡表》，两篇文章都是写给曹操的，这说明他曾经对曹操抱有很大的希望。他的书信和章表曾受到李充《翰林论》和刘勰《文心雕龙·章表》篇的称赞。我们只选了《论盛孝章书》一篇。陈琳的《为袁绍檄豫州》是一篇讨伐曹操的檄文，它以先声夺人之势，历数了曹操的种种罪状，并连及曹操的父祖，申明讨曹的目的是防止曹操篡权，以匡扶汉室。并号召各州郡起兵响应，在非常之时，建非常之功。文笔酣畅，气势若江河奔泻，是檄文中的优秀之作，对唐代骆宾王的《讨武曌檄》有明显的影响。

汉文的成就是相当高的，从唐宋的古文运动所宗法的文统来看，两汉文是它们的主要宗法对象。

韩愈学古文是从三代而至两汉，苏轼是“将以追两汉之馀，而渐复三代之故。”他们二人同样在提倡古文，韩愈主张由三代而下至两汉，苏轼则主张先追两汉，再逐渐地复三代之古，认为学习两汉文更是当务之急。到了明代的“前后七子”便提倡“文必秦汉，诗必盛唐”了。这从一个方面说明两汉文是古文发展的一个高峰，也是古文的始盛期。

东汉末的建安（196—220）时期，是一个社会大动乱的时期，也是“文学自觉”的开始，所谓“自觉”就是自我意识的觉醒，就是在文学中发现自我、表现自我，这又与思想解放相联系。这种新的风气也自然会影响到文的创作，文的抒情性加强了，对文的艺术形式美也更加讲究，曹丕在《典论·论文》中说“诗赋欲丽”，正是对文学作品提出的新的审美要求的表现。曹氏父子与“建安七子”的文章，文字都渐趋省净、清丽，在赋的领域，这种变化表现得更为明显，我们姑且以祢衡的《鹦鹉赋》和王粲的《登楼赋》为例，来看一看他们的赋有什么新的特点。

《鹦鹉赋》是咏物之作，但它意不在物，而是一篇托物言志之作。赋中所写的“挺自然之奇质”、“体金精之妙质”、“性辩慧而能言”、“才聪明而识机”，以及“容止闲暇，守植安停，逼之不惧，抚之不惊”等等，其咏物都带有象喻性，是其恃才傲物的性格写照。后文所写的“流飘万里，崎岖重阻”，“嗟禄命之衰薄，奚遭时之险巇”，更是他生不逢时、怀才不遇、为一个个统治者所不容而流徙播迁的写照，此赋可谓名为咏物实为咏怀。写鸟又写人，处处不离所写之物而又不粘于物，咏物与抒情融化无迹，辞采华丽而又朴茂清新，是难得之佳作。

《登楼赋》所抒写的，一为远适荆蛮的思乡之情，一为怀才不遇之叹。以他自己的身世遭遇而言，很有个性化的特点。它又采用了骚赋的形式，除了“兮”字之外，均为整齐的六字句，句式十分整齐，辞采明快华茂，音韵自然婉谐，全赋数次换韵，皆转换自如。在抒情手法上，采用了多种手法的交替，克服了大赋的平板，而显得富有波澜，体现了追求丽的倾向。它与骚赋有一脉相承的关系，而与大赋迥然有别。司马相如所奠定的大赋模式是散体，是将骚赋散文化，而汉末与曹魏时代的赋走的是抒情化与诗化的道路，是向骚赋的回归，但又不完全像骚赋。看来文学的发展也是沿着否定之否定的规律发展前进的。

建安至魏初抒情赋的代表作家是曹植，从现存资料看，他大约写了四十篇左右的赋，其中最引人注意的是赋苑中的奇葩《洛神赋》。作者以神奇的色彩，举体华美的文笔，绚丽的辞藻，丰富的想象，缠绵悱恻的感情，为读者塑造了一位光彩夺目的美神。此赋的写作时间比《登楼赋》晚十七年，在这十七年中曹植又将抒情赋推向了一个新的高峰。它明显地受到宋玉《神女赋》的影响，而不是汉代大赋的路子。它虽然使用了主客问答的形式，但这种问答既有别于宋玉的《高唐赋》与《神女赋》，也不同于司马相如的《子虚》、《上林》。《神女赋》反复使用对话，《子虚赋》等的对话带有辩难性质，而《洛神赋》的对话只是起了引起下文的作用，要而不繁，干净利落。从句式上看，《洛神赋》从三字句到十字句不等，长短句交错使用，纡徐委曲，错落有致。

曹植把赋变成美文，而且有诗化的倾向。

曹植的应用文写得也很美，我们所选的《求自试表》、《与杨德祖书》是其应用文的代表作。《求自试表》中已有些有意结撰的对偶句，为日后骈文的出现，奠定了初步的基础。

三国时代的文，蜀汉的文章，我们只选了脍炙人口的诸葛亮的《出师表》，曹魏的作家我们选了阮籍、嵇康和向秀的几篇文章。他们都是“竹林七贤”中的人物。曹魏时代是玄学盛行的时代，文章大受玄风的影响，玄学的思维是思辨性的，互相之间往往有所辩难，在这种风气的影响下，产生了一批论文。阮籍的《通易论》、《通老论》、《达庄论》、《乐论》等，都是玄学方面的论文，《大人先生传》也写得颇为精彩，这五篇文章因篇幅较长，我们没有选录，仅选了一篇奏记《诣蒋公》和一篇《猕猴赋》。阮籍在文中虽然没有说清楚不愿应太尉蒋济征召的原因，在天下多故的时代，为了远祸自全，他不愿裹进政治斗争的旋涡中去，只好言不由衷得敷衍几句，所以文章特短。阮籍是崇尚自然而反对名教的，对于虚伪的礼法之士他深恶痛绝，《猕猴赋》可视为讽刺礼法之士的小品。他以老、庄的崇尚自然本性的哲学观点，否认统治阶级历来宣传的禹铸九鼎“异物来臻”的祥瑞，又用猕猴来象喻追求功名利禄的礼法之士，把他们形容为人面兽心之徒，虚伪而善于表演，虽处缧绁，尚且作态媚人，一心希冀受到主人的嘉惠，终不免成为主人的玩物，最后落了个“伏死于堂下，灭没乎形神”的可悲下场。此赋构思奇特，描写惟妙惟肖，语言也生动活泼。

嵇康的文，我们选了《与山巨源绝交书》和《琴赋》两篇，这两篇文章都很有特色。嵇康是精通音乐的人，《琴赋》与《声无哀乐论》是他描写音乐、表现音乐美学思想的两篇重要文章。《琴赋》的同题之作不少，东汉的傅毅、马融、蔡邕都写过《琴赋》，而《文选》独选嵇作，是有眼光的。此赋从琴的材质、制作、演奏、琴德等诸多方面，加以铺张、渲染。从弹琴与听琴两个方面阐发作者的音乐美学思想，洋洋洒洒，仪态万方。其中的音乐描写，不仅使用了大量的形象比喻，而且涉及以视觉写听觉的“通感”问题，对后代的音乐描写影响很大。《与山巨源绝交书》是一篇刚肠疾恶而又带有几分偏狭的个性色彩的书信，书中所写的“七不堪”、“二不可”，是与官场决裂的宣言，也是嵇康人格的自我表白，可从一个侧面看到魏晋名士的风度和心态。颇有认识价值。

魏代的后期，大权已经落到司马氏手中，司马氏为了篡夺曹魏的政权，大力排斥异己，实行政治上的高压政策，当时名士少有全者。“竹林七贤”中的嵇康与吕安，就死在他的刀下。我们所选的向秀的《思旧赋》，就是为纪念嵇、吕两位亡友而作。鲁迅先生在《为了忘却的记念》中说：“年青时读向子期《思旧赋》，很怪他为什么只有寥寥的几行，刚开了头却又煞了尾。”在司马氏的高压政策下，有话不敢直说，所以便形成了文尚曲隐的现象。

晋代历史上分为西晋和东晋两个历史阶段。西晋的文和赋，与前代相比，在体式上没有太大的变化，只不过在语言上更追求藻饰。刘勰《文心雕龙·明诗》篇说：“晋世群才，稍入轻绮，采缛于正始，而力柔于建安。”这虽然是针对诗歌的发展而言，但也部分地适用于文的情况。这一时期的文，语言渐趋骈偶化。就文而言，是由散体的文在逐渐向骈体文过渡；就赋而言，处于由辞赋向骈赋的过渡时期。西晋的文，我们选了李密的《陈情事表》、张华的《鹪鹩赋》、左思的《三都赋》与《白发赋》、潘岳的《秋兴赋》与《闲居赋》、陆机的《叹逝赋》等，以上这些作品除《白发赋》外都见于《文选》。同时我们也选录了四五篇《文选》所不曾注意的作品，如刘琨的《壶关上表》、《与石勒书》、《答卢谌书》和鲁褒的《钱神论》等，以增加应用文和其他文体的篇数。李密的《陈情事表》由于《文选》和《古文观止》都选了它，故最为脍炙人口，表现作者与其祖母相依为命的亲情，尤为感人肺腑。张华的《鹪鹩赋》是托物言志之作，与祢衡《鹦鹉赋》的手法是一脉相承的。左思的《三都赋》问世之后，一时争相传写，有洛阳纸贵之说。这篇写京都的大赋，走的还是汉代大赋的路子。汉代大赋描写京都题材的有四五篇之多，如班固的《两都赋》，张衡的《西京赋》、《东京赋》、《南都赋》等，均为《文选》所选录，而且所选作品的数量，在赋的十五种分类中居于首位。而这种题材的赋由汉至晋经久不衰，仍有人在蹈袭前人之路在仿效摹写。左思在完成《齐都赋》的创作之后便开始准备，从资料的收集、构思到写成，他大约用了十余年的时间，是他的力作。刘勰《文心雕龙·才略》篇云：“左思奇才，业深覃思，尽锐于《三都》，拔萃于《咏史》。”就是把《三都赋》作为左思文的代表作的。另一篇《白发赋》是用拟人化的手法，托白发与作者的对话而写成的寓言小赋，这篇赋可能作于晚年，他把一生未得志的牢骚诉之于诙谐而幽默的对话之中，显得姿趣横生，新颖别致。潘岳的两篇赋，不管是写景还是抒情，都能达到情景交融的境界，辞藻绚丽明净，行文自然流畅，自有一种清丽的风韵。

刘琨的文大多写于后期，即永嘉元年（307）任并州刺史之后，我们所选的三篇应用文，都很有特色。西晋后期，内乱外患交相发生，历时十余年的“八王之乱”，使中原大地的人民蒙受了巨大的灾难，这时居于我国北部的少数民族的首领们，也乘虚而入，西晋政权受到严重的威胁，刘琨就是在这种社会背景下任并州刺史而走上抗敌前线的。在上任的途中，他所看到的现实是：“目睹困乏，流移四散，十不存二，携老扶弱，不绝于路。及其在州者，鬻卖妻子，生相捐弃，死亡委危，白骨横野，哀哭之声，感伤和气。群胡数万，周匝四山，动足遇掠，开目睹寇。”（《为并州刺史到壶关上表》）其描写当时的现实是非常真实可信的。在西晋末，刘琨官至大将军、司空、都督并、冀、幽三州诸军事，是当时的北方重镇。为了匡复晋室，靖乱中原，他在抗敌前线做了许多工作。他的《与石勒书》是为分化瓦解敌人劝石勒降晋的，《书》中对石勒晓之以大义，说之以利害，文字简净，结构顿挫起伏，文笔清丽，一气贯注，在晋文中是不可多得的。《答卢谌书》是刘琨遇害前的作品，《书》中除了抒写与卢谌分别的怅恨与勉励卢谌之外，还说到自己的思想转变过程，在国家危机之秋，在现实的教育与感召激励下，促使刘琨从老、庄思想中觉醒过来。从魏晋名士追求的放纵、旷达中解脱出来，他说：“昔在少壮，未曾检括，远慕老庄之齐物，近嘉阮生之放旷。怪厚薄何从而生，哀乐何由而至。自顷辀张，困于逆乱，国破家亡，亲友凋残。负杖行吟，则百忧俱至；块然独坐，则哀愤两集。……然后知聃、周之为虚诞，嗣宗之为妄作也。”这就是刘琨觉醒后的自述。此文颇有认识价值。文章言简意赅，骈散结合，可见晋文渐趋骈偶的发展趋向。

鲁褒的《钱神论》是过去不曾注意的一篇闪耀着唯物主义思想光辉的文章，它用冷嘲热讽的语言，淋漓尽致地揭露出钱能通神的现象。作者对金钱拜物教的批判是与黑暗腐朽的西晋社会联系在一起的，这比沙氏比亚在戏剧中对金钱拜物教的批判要早一千多年。作者把孔圣人所说的“生死有命，富贵在天”改为“生死无命，富贵在钱”，这是非常大胆的，既批判了儒家的“天命论”，又对金钱的经济杠杆作用，作了充分的估计，其经济思想达到了相当的高度。作者对货币在流通过程中的作用，并未否定，所否定的是统治阶级对金钱的贪婪与无限的追求。此文很富有幽默感，把铜钱称为“孔方兄”是从此文开始的。

东晋士族文人重玄谈的风气与西晋大体相似，此时佛教的流行，又更盛于西晋。偏安江南的东晋小朝廷，一代代的皇帝大多不思振拔，不图恢复，过着苟且偷安的生活。这一时期的文学，在文和赋的领域，渐多模山范水之作。郭璞的《江赋》、孙绰的《游天台山赋》可作这一方面的代表。表现祖国的山水之美，为文艺园地增添新的花朵，本来是一件好事，遗憾的是他们对广阔的社会现实反映得太少，已少有像刘琨那样的作家了。东晋的文值得称道的是王羲之的《兰亭集序》和陶渊明的几篇文章。《兰亭集序》融记叙、议论、抒情于一体，语言朴素自然，对后代散文影响颇大。序中还写到作者自己的人生觉醒，他经过人生的反思，认识到“一死生为虚诞，齐彭殇为妄作”。对庄子的“齐物论”提出批判，在老庄哲学与玄学风行的时代，作者有此认识是难能可贵的。这是刘琨之后的又一个觉醒者。

东晋的文以陶渊明的成就为最高，这与他的“质性自然”与恬淡的心境有关。《归去来兮辞》是其代表作。这是一篇陶渊明与官场决裂的宣言书，他率真地向读者打开了他的心扉，他把对官场的认识，对人生真谛的思索，对做官与弃官归田的原因，对前半生人生道路的反思与觉醒等等，都和盘托出。他的心态，他的精神境界，均在一篇《归去来兮辞》中得到充分的表现。它在艺术上也有独到之处，语言清丽、简净，自然而又有韵味，千百年来赢得许多人的称赞。欧阳修说：“晋无文章，唯陶渊明《归去来辞》而已。”这话说得虽然有点绝对，但这篇文章确实是晋文中的佼佼者。《桃花源记》是一篇闪耀着理想光辉的优美散文，桃花源是作者理想中的天国，是没有压迫、没有剥削、没有尔虞外诈，民风淳朴的一片净土，它与污浊的官场和黑暗的社会现实是对立的，而且有深厚的思想渊源。它在艺术构思上颇为巧妙，奇中求真，引人入胜。且笔法疏淡，语言简约，在时代文风渐趋骈偶的东晋，陶渊明独能拔出流俗，采用纯粹的散文体来写作。他的《五柳先生传》其语言风格也大体如此。而其《闲情赋》确是别具一格的文章，赋中描写了一位令人日夜悬想的绝色佳人，作者幻想日夜与她相处，形影不离，甚至想变成各种器物，附着在这位美人身上。这完全是人们心目中的另一个陶渊明，其主题思想历来就有不同的说法，为了避免穿凿附会，我们比较倾向于爱情说。

南朝宋的文向骈体的发展又进了一步，刘勰说：“宋初文咏，体有因革，老庄告退，而山水方滋。俪采百字之偶，争价一句之奇，情必极貌以写物，辞必穷力而追新：此近世之所竞也。”（《文心雕龙·明诗》）这虽然是针对诗歌而说的，也基本上适用于文，这时山水题材的文多了起来。我们所选的谢惠连的《雪赋》和谢庄的《月赋》，就是这方面的代表作。《雪赋》用赋的传统的主客对话的形式，假托西汉的梁孝王在兔园赏雪，召来了枚乘，邹阳、司马相如三人一起咏雪，主客各借典故，刻画雪景。司马相如的一大段咏雪是赋的主体，写雪景之美，瑰丽奇绝，极尽化工之妙。赋中又有枚乘的“积雪之歌”和“白雪之歌”，添加了携佳人饮酒赏雪的艺术描写，增强了赋的抒情成分。枚乘的乱曰，由赋景、抒情而转向议论，赞美了雪的品格，将雪人格化。雪的随物赋形、从风飘零、遇热而融、遇染而变、无虑无营的特点，正是老庄随遇而安、委运任化、恬淡自然思想的反映。这种结尾与当时山水诗的玄言尾巴有点相似。《月赋》也采用了“假主客以为辞”的传统赋法，用曹植等四位文学家作为赋中人物，“先叙事，次咏景，次咏题，而终之以歌，从首至尾，全用《雪赋》之格。”（祝尧《古赋辨体》卷六）比起《雪赋》来，《月赋》的议论和玄言尾巴都不见了，它写月之美，运用的是烘云托月之法，重在借月色抒情，抒写出在深秋季节的寂寞月夜中，人们对月的感受。赋末的歌“美人迈兮音尘阙，隔千里兮共明月”云云，又着上一层月夜怀人的凄楚气氛。从《雪赋》到《月赋》，我们可以看出宋初赋的一些演进之迹。《雪赋》和《月赋》在形式上已基本上可说是骈赋了。骈赋也称俳赋、“四六赋”，它是汉赋（辞赋）的变体，在体式上的特点是重对偶，句式多为四字句和六字句，四六交错使用而形成对偶，这就是所谓的“骈四俪六”。在用辞上讲究辞藻绮丽，同时也讲究声韵的婉谐。读之声韵铿锵，朗朗上口。其后鲍照《芜城赋》的出现，标志着骈赋的发展已臻于成熟的境地。

《芜城赋》在体式上与《雪赋》、《月赋》不同的一点是“假主客以为辞”的汉赋模式完全脱尽，既无人物也无对话，纯粹是抒情主人公的语言，句式也更加整齐，可视为骈赋的代表作。此赋将广陵昔日的繁盛与今日的萧条进行对照描写，写今与昔都使用了艺术的夸张，写昔日之盛正是为了突出今日之衰，读之让人惊心动魄。它在语言的提炼上，于浓郁之中见奇峭，骈俪之中有转折、错落、变化，所以获得人们的交口称赞。

鲍照在宋代是位文学创作成就很高的作家，因出身寒门，受世族门阀的压制，在仕途上很不得志。“才秀人微”的鲍照在政治上是有远大抱负的人，理想与现实的矛盾郁积在他的胸中，遇之即发，这在我们所选的《登大雷岸与妹书》和《瓜步山楬文》中都有所表现。《登大雷岸与妹书》通篇都是写景，在前此的书信中是不曾看到的，流露出借山水写怀抱的倾向。“窥地门之绝景，望天际之孤云，长图大念，隐心者久矣。南则积山万状，负气争高，含霞饮景，参差代雄。”那隐心已久的“长图大念”与“负气争高”的精神，正是鲍照此时心境的写照。描写庐山的一段，色彩绚丽，真乃画笔难到。这是一篇优秀的骈文书信。

《瓜步山楬文》由登高望远引出“意类交横”的感慨，又从“数寸之钥”与“千钧之关”的关系，引发出人的“才施”与所处地位的关系。接着笔锋一转，说瓜步山的一切特点完全是“居势使之然也”，于是发出“才之多少，不如势之多少远矣”的感慨。真乃点睛之笔。这是作者对门阀制度的抨击，是才智之士受压抑的一声吼叫，在当时的文坛上这种批判现实的呼喊是很少见的，它如孤星夜月，发出熠熠的光彩。后文又抨击了“贩交买名”之辈，“吮痈舐痔”之徒，他们的是是非非，是不屑一评的。此文感情激烈，笔力遒劲，气势雄放，大有睥睨千古之慨，是不可多得的佳作。

齐梁时代的文，骈体已占绝对的优势，赋则多为骈体赋，文则多为骈体文，就连应用文也受到这种风气的影响，日趋骈偶化。而且“永明声律说”的出现，沈约等人为文皆用宫商，将平上去入的四声理论用之于诗歌创作中，这就促使骈文在声韵上运用平仄，以使声韵和谐。又由于类书的出现，隶事用典的风气也盛行起来，六朝的骈文就这样形成了。骈文在当时是被目为“美文”的，一时蔚成风气，章、表、檄、移，诏、令、碑、铭，乃至书信等，无不受其影响，出现了大量的骈体之作。但骈文的名称出现较晚，清代的李兆洛说：“自秦迄隋，其体递变，而文无异名。自唐以来，始有古文之目，而目六朝之文为骈俪，其为学者，亦自以为与古文殊格。”（《骈体文钞·序》）从唐代开始，古文与骈文成了两个不同的派别，此后不断有人反对骈文，形成骈散之争，但直到清末，骈文仍在流行，这也说明它是有生命力的。骈文充分利用了汉字的声律特点和可对偶性，它的排比对偶的修辞手法以及对形式技巧的讲求，还是值得继承和借鉴的。

永明作家，我们选了沈约的《愍衰草赋》和谢朓的《拜中军记室辞随王笺》，两篇文章都是四六文，《愍衰草赋》诗化的倾向十分明显，后半篇的“风急崤道难，秋至客衣单。既伤檐下菊，复悲池上兰。飘落随风尽，方知岁早寒”云云，直至篇末，纯为整齐的五言诗句。谢朓的《拜中军记室辞随王笺》，是典型的四六文，如：“不寤沧溟未远，波臣自荡；渤澥方春，旅翮先谢。清切藩房，寂寥旧荜；轻舟反溯，吊影独留。白云在天，龙门不见，去德滋永，思德滋深。唯待青江可望，候归艎于春渚；朱邸方开，效蓬心于秋实。”其对偶是颇为工整的。

江淹的《恨》《别》二赋，是脍炙人口的骈赋佳作，他写出了各种不同的生死之恨与“别则一绪，事乃万族”的离情别恨，这种永恒的主题，很容易引起人们的共鸣。特别是《别赋》，既写出了离别之恨的共性，也写出了它的个性，不同层次的人物各有不同的特色：有的缠绵悱恻，有的慷慨悲壮，有的哀婉凄绝，有的铭心刻骨。又加作者善于用自然景物来渲染离别的气氛，以景衬情，使离恨倍增其凄凉色彩。写游子之别，转笔写“居人愁卧”的种种情景，两相映衬，相得益彰。“春草碧色，春水渌波，送君南浦，伤如之何！”是千古传诵的名句。

孔稚圭的《北山移文》也是一篇骈文，它淋漓痛快地讽刺揭露了“身在江湖，心存魏阙”的假隐士，在艺术上获得很高的成就。许梿评它说：“此六朝中极雕绘之作，炼格炼词，语语精辟，其妙处尤在旋转得法，当与徐孝穆《玉台新咏序》，并为唐人轨范。”又说它“造语精缛，却无一字拾人牙慧”（《六朝文絜》卷八）。

丘迟的《与陈伯之书》和吴均的《与宋元思书》，都是当时书信中的名篇，它们的出名，都与书信中的景物描写有关。《与陈伯之书》中的“暮春三月，江南草长，杂花生树，群莺乱飞”是写景的名句。《与宋元思书》通篇都是写景，这种先例是鲍照的《登大雷岸与妹书》开创的。但鲍照的《登大雷岸与妹书》借山川形胜写出了自己的“长图大念”与“负气争高”的胸襟，这一点又非吴均所可追攀。

梁代的文人多在萧氏父子的几个文学集团之中，萧氏父子的文学观点不尽相同，以萧统、萧纲、萧绎三兄弟而论，萧统的文学观儒家传统的东西比较多，而萧纲、萧绎则属于新变派。萧氏三兄弟的文章，我们选了萧统的《陶渊明集序》，萧纲的一篇书信和萧绎的两篇小赋。《陶渊明集序》还是骈散结合的，文中既阐述了作者自己的人生哲学思想，也概括了陶渊明的创作旨趣，高度评价了陶渊明的创作，一语中的地指出陶渊明诗歌的一大特点：“吾观其意不在酒，亦寄酒为迹焉。”对陶渊明作品的艺术成就也给予了高度的评价，唯对《闲情赋》的评价，有些迂见，所以遭到后人的非议。萧绎的《采莲赋》与《荡妇秋思赋》颇有特点，篇幅短小精悍，文字清新流丽，是诗化的骈赋。庾信前期的赋，如《荡子赋》等，也是多以整齐的诗语入赋，语言风格大体与萧氏兄弟相近。

南朝陈的文和赋，我们所选不多，仅选了徐陵的《玉台新咏序》和张正见的《石赋》两篇作品。《玉台新咏序》用虚拟的手法，假托《玉台新咏》之撰集，是出自“本号娇娥，曾名巧笑”的丽人之手，这种写序的手法，可谓别出心裁。序中只强调诗歌的审美愉悦作用，而只字不提文学的教化作用，流露出鲜明的新变派观点。此序在艺术上颇为引人注意，许梿评论它说：“态冶思柔，香浓骨艳。”“骈语至徐、庾，五色相宣，八音迭奏，可谓六朝之渤澥，唐代之津梁。而是篇尤为声偶兼到之作。炼格炼词，绮绾绣错。几于赤城千里霞矣。”（《六朝文絜》卷八）

张正见的《石赋》是一篇咏物的事类赋，赋虽不长，但用典甚多，每个典故都与石有关，其中有不少关于石的神话传说，为此赋披上了一层神奇的色彩。以石名篇的赋此篇最早，它对唐以后《石赋》的写作有一定的影响，在赋史上应占有一席地位。它的用典很有特点，它把许多典故先进行艺术化的加工，再组织在对偶化的骈俪文字中，花费了极大的艺术匠心。李调元说：“陈张正见《石赋》，如‘鱼跃湘乡之水，雁浮平固之湖’，‘堕山鹊之金印，碎骊龙之宝珠’，通章无句不对，实开律赋之先。”（《赋话》卷一）

北朝的文，作家之众，作品之多，虽然比不上南朝，但从文风上看，也自有它不同于南朝的清刚之气，又加上庾信、颜之推等作家都是由南入北的，这也为北朝一向比较沉寂的文坛增添了几分光彩。

北朝有三部著名的用散文写成的著作，即郦道元的《水经注》、杨衒之的《洛阳伽蓝记》和颜之推的《颜氏家训》，我们各选了个别的篇章、段落，以见一斑。其中《水经注·江水注》有关山峡一段的描写，最为脍炙人口。

北朝的赋，我们选了卢元明的《剧鼠赋》，这是一篇含有讽刺寓意的咏物小赋，赋中对剧鼠的丑态做了惟妙惟肖的描绘，又运用了我们古代有关鼠类的许多典故，旨在以鼠影射危国群小，其中包括老奸巨猾的贪官污吏和善于朋比党奸的溜须拍马之徒，意蕴含蓄而丰富。此赋在语言上偶句颇多，也属于骈赋。

北魏的温子昇是位颇有才能的作家，据说庾信初次出使北方，非常看不起北方的文士，唯对温子昇的《韩陵山寺碑》读后抄写下来，认为其馀的作家作品都不屑一顾。（见张鷟《朝野佥载》卷六）这篇碑文把尔朱兆叛乱的时代背景和碑主高欢的丰功伟绩，全用典故表现，用典也十分贴切，行文多用对偶句，而且很有文采，是典型的骈体文。与南朝的骈文初无二致，难怪为庾信所欣赏。

庾信是集六朝大成的作家，他的文以他四十二岁羁留北朝为界，分为前后两个时期。上文所提到的《荡子赋》，可作为前期文的代表。由于当时的生活视野比较狭窄，与萧纲、萧绎的赋并无多大的不同。入北之后，由于生活经历的变化，仕北的惭耻与“乡关之思”时时撞击着他的心灵，故能以意类纵横的凌云健笔，写出了“暮年诗赋动江关”的杰作。《小园赋》和《哀江南赋》，便是这方面的代表作。庾信在北周没有隐居过，他所刻画的小园不过是庾信心造的泡影、理想的天国，后半篇以乡关之思转为哀怨之辞，将内心的痛苦写得相当深刻。为了求得精神上的解脱，他只好乞灵于老庄。此赋受陶渊明《归去来兮辞》和《桃花源记》的影响比较明显，他心目中的小园，实为躲避人世风雨的世外桃源。似乎参考了陶渊明的蓝图。在赋的发展史上，它是六朝骈赋向唐代的律赋演变过渡的产物。李调元曾经指出：“古变为律，兆于吴均、沈约诸人。庾子山信衍为长篇，益加工整。如《三月三日华林园马射赋》及《小园赋》，皆律赋之所自出。”（《赋话》卷一）但代表庾信赋的最高成就还是他的《哀江南赋》，它是骈体赋的鸿篇巨制，具有史诗的规模和气魄，是一篇划时代的杰作，在赋史上具有里程碑的意义。它的内容相当丰富和深厚，是对梁朝兴亡的历史总结。它艺术地再现了太清二年（549）的侯景之乱和承圣三年（554）的江陵之乱，多层次、多侧面地描绘出两次战乱中的种种生动画面，涉及众多人物在动乱中的表现。既歌颂了卫国的英雄，也鞭挞了梁宗室子弟的坐观时变和卖国求荣之辈，批判了梁武帝的沉溺佛讲和不修武备，以及梁元帝的猜忌凶残，同时还反映了被虏至长安的十万俘虏的血泪生活。基于此，所以我们说它具有史诗的规模的气魄。

《哀江南赋》又是自传体的史诗，它几乎概括了庾信一生的经历，并涉及他的列祖列宗。它将个人的经历与家国的历史，即将自传与史诗结合起来描写，十分突出地表现了自己的“乡关之思”与亡国之痛。

《哀江南赋》的艺术成就很高。在语言方面，它大量使用四六对句，于对偶之中，时杂散体语言，改变了他前期赋作以大量诗语入赋的倾向，富有错综变化之美。在声律方面，宫商抑扬，声韵婉谐。在文采方面绣错绮交，形象生动，光色鲜艳。其隶事用典尤为突出，它把一个个互相没有联系的历史典故，连成完整而又有象喻性的意象，“援古证今，用人若己”（《文心雕龙·事类》）。总之，庾信在驾驭骈赋的形式技巧方面，已经达到了炉火纯青的境地。庾信在赋史上的地位，可以说是集六朝之大成，启唐之先鞭。

以上通过一些有代表性作品的分析评价，粗略地勾勒出汉魏六朝文的发展演变及艺术成就。本书精选50家72篇，依丛书体例，文前撰有作家简介，文后有详细注释，注〔1〕交代所用版本，并对文章稍加点评，希望对读者的阅读有所帮助。责编葛云波先生在本书修改过程中提出不少宝贵的意见，并纠正了不少错误，在此谨表谢忱。书中不当之处，还请广大读者和海内外专家指正。

\begin{flushright}
刘文忠\\
2002年12月初稿\\
2009年12月改定
\end{flushright}